\chapter{实验目的}\label{part:purpose}

在日常课程学习与编写代码的实践里，我们小组成员频繁使用大语言模型辅助推导数学题、检索长篇英文文献以及调试 Python 脚本。随着使用频次增加，几个反复出现的现象引起了我们的注意：同一道数学题，有时要求模型把步骤写清楚能大幅提高正确率，有时却因为推导写得过于冗长繁琐而在输出上限处被硬生生截断；把二十篇参考资料一起塞进对话框时，模型经常漏掉夹在中间的关键句子；而在让模型修改报错代码时，单纯催促它“再仔细检查一遍”，往往不仅修不好原有漏洞，还会把原本已经写对的逻辑改坏。

面对这些零散的日常经验，《数据分析与可视化实践》课程设计题目二“AI 对话数据与智能体行为分析”为我们提供了一次用严格数据说话的机会。我们小组选择\textbf{选项 B：构建可视交互系统}，围绕上述实际痛点设计并实现了名为\textbf{“知微”}的多任务模型行为可视分析系统。之所以取名“知微”，寓意从单次 Token 消耗、单个证据段落位置和单轮报错反馈这些细微日志特征入手，观察并验证模型整体表现的变化规律。

\section{课程设计任务定位与总体目标}

对照课程设计指导书对“可视交互系统组”的考核要求，我们没有停留在把现成对话记录做成几张静态柱状图的层面，而是将工作重心放在“日志采集与脱敏清洗—同题配对统计推断—多视图联动探索系统”三个相互咬合的环节上。图~\ref{fig:roadmap} 概括了小组从最初拟定研究问题到最终完成无头浏览器回归验收的六个实施阶段。

\begin{figure}[!htbp]
\centering
\includegraphics[width=0.95\linewidth]{\detokenize{figures/diagrams/d01-roadmap.pdf}}
\caption{课程设计六阶段推进路线与阶段产出物关系}
\label{fig:roadmap}
\end{figure}

在具体实施过程中，我们为系统确立了以下四项核心实验目标：

\paragraph{规范化采集与隐私伦理保护}针对多模型、多提示策略与多轮代码执行日志建立完整的原始备份机制。所有原始 JSONL 响应文件在落盘后即做只读冻结与 SHA-256 摘要登记，一切格式转换、评分修正与字段抽取均在派生副本上完成。在对外展示与仓库交付层，严格剥离 API 密钥、本地绝对路径与账户信息，并对原始题目编号执行实验域内的单向哈希映射（\texttt{question\_key}），做到既保留同题配对分析能力，又不公开未授权的原始长对话文本。

\paragraph{异构日志整合与空值语义隔离}本项目的实验横跨数学与多选推理、多文档问答、单轮多采样代码生成以及多轮交互代码修复四类不同场景。不同场景天然带有各自的评测标尺与专有字段。我们在清洗阶段制定了 21 列统一匿名特征规范（\texttt{all\_records.csv.gz}，详细字段字典供查阅见附录~\ref{app:db}），统一布尔值、UTC 时间戳与截断判定口径，并严格区分“正常计零分（\texttt{0}）”“沙箱环境不兼容排除（\texttt{environment\_excluded}）”与“源端未记录（\texttt{null}）”，避免用草率的默认值填补污染分母。

\paragraph{尊重题目相关性的统计建模}由于同一道基准题会被多个模型、31 种提示词策略乃至 5 次重复采样反复作答，单条响应之间存在明显的题目聚类效应。如果直接把二十多万条记录当作彼此独立的伯努利试验计算标准误，置信区间会被人为压窄。因此，我们在分析阶段全面采用\emph{以独立题目为聚类单元的同题配对 Bootstrap 重采样}（默认重采样 $B=5{,}000$ 次，固定随机种子 91），并且严守描述性统计与因果解释的界限，不把特定测试集上的相关性夸大为普适定律。

\paragraph{多层次可视交互与证据溯源}为了兼顾首次浏览时的结论可读性与深入排查时的自由探索需求，我们构建了双模式前端体系：主入口（\texttt{index.html}）仅依赖浏览器原生 HTML5、CSS3 与手写 SVG 构建五球导航与五屏联动视图，老师或同学克隆仓库后只需用 Python 3.10 标准库即可一键启动；探索工作台（\texttt{explore.html} 与 \texttt{legacy.html}）则离线集成了课程指导书推荐的 ECharts、D3.js、Plotly.js 与 Vega-Lite 四类工具，支持二维散点框选、采样步进播放、分布缩放以及单题代码 Diff 抽屉溯源。

\section{核心研究问题拆解与假设边界}

\Needspace{15\baselineskip}
\begin{wrapfigure}{r}{0.47\linewidth}
\vspace{-8pt}
\centering
\includegraphics[width=\linewidth]{\detokenize{figures/diagrams/d08-research-cube.pdf}}
\caption{三维实验操控空间与五项发现映射}
\label{fig:research-cube}
\vspace{-8pt}
\end{wrapfigure}

如果把大模型本身看作参数已经固定的黑盒求解器，普通使用者在实际调用接口时真正能够主动调节的抓手主要有三个：\textbf{怎么提输出要求}（提示词组织与长度约束）、\textbf{怎么摆参考材料}（关键证据在长上下文中的位置），以及\textbf{怎么给纠错信号}（多轮交互中的测试反馈粒度）。如图~\ref{fig:research-cube} 所示，我们围绕“\emph{在不改变模型底层权重的前提下，如何组织输出表达、参考证据与执行反馈以兼顾任务质量与计算开销？}”这一总问题，设计了五项具体、可证伪的子问题，恰好对应课程指导书要求的两条低维模式、两条三维及以上交互关系和一条反常现象：

\begin{enumerate}
  \item \textbf{低维问题一（输出组织 $\to$ 质量与开销折衷）}：在数学推理任务中，从只给最终答案（\texttt{NoCoT}）切换到分步推理（\texttt{DefaultCoT}）通常会伴随数倍的 Token 增长。如果在保留分步思考的前提下加上一句温和的简洁提示（\texttt{BeConcise}），能否在数学、问答与代码任务上普遍省下可观的输出长度而不牺牲准确率？
  \item \textbf{低维问题二（作答语言模板 $\to$ 竞赛数学截断与得分）}：面对同一批英文原版数学题（如 \texttt{MATH500} 与 \texttt{GSM8K}），仅将外显作答要求改为中文（\texttt{ChineseCoT}），模型输出 Token 数、达到上限被截断的比例以及最终题目得分会发生怎样的联动变化？这种变化在不同难度题库之间是否存在明确边界？
  \item \textbf{反常现象问题（极端长度软约束 $\to$ 输出反弹与格式崩溃）}：直觉上，把提示词里的字数、词数或字母数上限压得越低，模型吐出的文本就应该越短。然而当约束收紧到“只准用 1 个词”或“只准用 10 个字母”这类极端阈值时，实际输出 Token 数是否真的单调下降？为何苛刻的格式指令会在部分模型上诱发严重的超长输出反弹甚至代码提取失败？
  \item \textbf{多维交互问题一（受测模型 $\times$ 证据位置 $\times$ 文档总数 $\to$ 问答命中率）}：在包含 10 篇或 20 篇检索文档的开放域问答中，保持问题与候选段落内容完全不变，仅把唯一包含标准答案的文档依次放到第 1 至第 20 位的不同位置上，不同模型表现出的首尾偏好（U 型曲线）深度有何差异？两代同系列模型之间的准确率优势在首段、中段和末段是否均匀分布？
  \item \textbf{多维交互问题二（受测模型 $\times$ 反馈信息粒度 $\times$ 修改轮次 $\to$ 代码强测通过率）}：当模型初稿代码未能通过测试时，连续 4 轮给予不同粒度的外部反馈——从完全不给反馈的独立重采样、无外部信号的自行检查，到仅告知通过与否、提供报错行定位诊断，再到给出单个或多个具体失败反例——模型的累积修复率、把原本正确初稿改坏的回退率以及总 Token 账单呈现怎样的三维演化规律？
\end{enumerate}

\CaseCard{方法论边界先行约定：相关性不等于内部机制因果}{
在进入后续实验分析之前，我们小组内部统一了一条贯穿全报告的写作准则：本实验采集到的全部指标均来自接口外显输入与返回文本的对照观测。例如“中文作答模板在 \texttt{MATH500} 上得分更高”，只能严格解释为\textbf{在 800 Token 输出上限条件下，中文表层表述密度更高、减少了未完结截断失败}，而绝不能凭此断言模型底层神经元“用中文思考更聪明”；同理，各提示策略在特定题集上的排序属于探索性统计证据，不能脱离采样温度与评测规则泛化为绝对优劣。}

\section{实验报告内容组织与正附分离约定}

由于整个课设周期内小组累计采集并归档了六批次共 273,186 条物理模型响应、编写了 158 个前后端与分析脚本文件（共 37,534 行代码）并产出了数十张分条件统计表，若将全部字段定义、提示词原文和分模型明细表直接堆入正文，不仅会打断分析逻辑，也会造成正文版面拥挤。

因此，本报告在编排上采取\textbf{“正文重机制推演与视觉证据、附录重规范备查与完整归档”}的结构设计。正文第~\ref{part:team} 章介绍四人分工协作机制与版本迭代过程；第~\ref{part:project} 章阐述“知微”系统的双口径数据全景、三入口界面形态与四层零依赖架构；第~\ref{part:process} 章详细展开实验设计、21 列清洗与评分器修正、五项核心统计发现、原生 SVG 与四库联动前端实现以及自动化质量回归；第~\ref{part:conclusions}、\ref{part:experience} 章总结实证启示、局限性与小组成员的真实课设收获。至于承担查阅功能的完整数据字典、31 种提示模板与模型规格、关键核心源码、原始问答与代码 Diff 案例库、多视口系统实测图集、全维度统计大表、工程目录映射及评分点核验日志，则按主题分类编排为\textbf{附录 A 至附录 H}，方便随时按需翻阅核查。
